\section{The Great Fez Hat Sanction}

Sanctions are nothing new. During the years before World War I, the Austria-Hungary Empire was a great power and manufacturing powerhouse. The Czech Republic, the part of the Empire, made most of the Fez hats in the world. The Fez was important to the Ottoman Empire:

\begin{quotex}
The fez became a symbol of the Ottoman Empire in the early 19th century. In 1827, Mahmud II mandated the fez as a modern headdress for his new army, the Asakir-i Mansure-i Muhammediye. The decision was inspired by the Ottoman naval command, who had previously returned from the Maghreb having embraced the style. In 1829, Mahmud issued new regulations mandating use of the fez by all civil and religious officials. The intention was to replace the turban, which acted as a marker of identity and so divided rather than unified the population.
\flright{\textit{Wikipedia}}
\end{quotex}


During the first decade of the 20th century, there were several wars between various Balkan nations and the Ottoman Empire. Meanwhile, Austria-Hungary had its own designs on the South Slavs. It had been gradually getting closer to Bosnia and Herzegovina. But when the Empire formally annexed Bosnia in 1908, Turkiye put sanctions on Austrian Fez hats. As we know, sanctions hurt the sanctioner, so Turkiye lost their hat supply. Turks were then forced to turn to other hats.

In any event, Ataturk banned them a decade later as part of the effort to modernize Turkiye.


\flrightit{Posted on 2022-10-16 by Cologero}

